% GENERATED FILE. Edit canonical lessons, then run npm run generate:books.
% canonical-source-hash: fnv1a64:ab7ccd67353dad04
% canonical-lessons: BN-C111-dori, BN-C111-perek, BN-C111-haturi, BN-C111-kurul, BN-C111-chhata

\chapter{Rope, Nail, Hammer, Axe, Umbrella}
\label{ch:bn-rope-nail-hammer-axe-umbrella}
\hlchaptermodality{111}

\begin{chapteropening}
\textbf{By the end of this chapter:} \emph{I can say a rope, a nail, a hammer, an axe and an umbrella.}

Complete the last lesson of chapter 111: I can say a rope, a nail, a hammer, an axe and an umbrella.
\end{chapteropening}

\section[\texorpdfstring{\bn{দড়ি} (dôṛi)}{দড়ি (dôṛi)}]{\bn{দড়ি} (dôṛi) — a rope}

\label{lesson:BN-C111-dori}



\begin{quote}
Before the new one: say the Bengali for matches, then the Bengali for coal.
\end{quote}

\subsection*{You'll want to know: \bn{দড়ি}}
\textbf{\bn{দড়ি}} — \emph{dôṛi} — \textquotedblleft{}a rope\textquotedblright{}.

\begin{cousinweb}
\textbf{\bn{দড়ি}} (\emph{dôṛi}) is a tadbhava, a word inherited from Sanskrit and reshaped by sound change. It comes from \emph{dṛti}, \textquotedblleft{}skin, leather bag\textquotedblright{} (CDIAL 6511). CDIAL glosses Bengali \emph{daṛi} \textquotedblleft{}sinew, tendon, cord, string\textquotedblright{}.
\end{cousinweb}

\begin{grammarlens}[title={What you've built}]
One more word for this chapter.
\end{grammarlens}

\subsection*{Your turn}
\begin{itemize}
\raggedright
  \item \emph{Say it:} \emph{dôṛi}
  \item \emph{Say it:} \emph{dôṛi}, once more
  \item \emph{Say it:} say \emph{dôṛi}
  \item \emph{Recall:} say the Bengali for matches, then the Bengali for coal, then say \emph{dôṛi} again
\end{itemize}

\begin{tcolorbox}[breakable,colback=teal!4,colframe=teal!35!black,title={Before you move on}]
What does \bn{দড়ি} mean? (\textquotedblleft{}A rope\textquotedblright{}.) Say it once more.
\end{tcolorbox}
\section[\texorpdfstring{\bn{পেরেক} (perek)}{পেরেক (perek)}]{\bn{পেরেক} (perek) — a nail}

\label{lesson:BN-C111-perek}



\begin{quote}
Before the new one: say the Bengali for coal, then the Bengali for a rope.
\end{quote}

\subsection*{You'll want to know: \bn{পেরেক}}
\textbf{\bn{পেরেক}} — \emph{perek} — \textquotedblleft{}a nail\textquotedblright{}.

\begin{cousinweb}
\textbf{\bn{পেরেক}} (\emph{perek}) is a loan from Portuguese \emph{prego}, \textquotedblleft{}nail\textquotedblright{}.
\end{cousinweb}

\begin{grammarlens}[title={What you've built}]
One more word for this chapter.
\end{grammarlens}

\subsection*{Your turn}
\begin{itemize}
\raggedright
  \item \emph{Say it:} \emph{perek}
  \item \emph{Say it:} \emph{perek}, once more
  \item \emph{Say it:} say \emph{perek}
  \item \emph{Recall:} say the Bengali for coal, then the Bengali for a rope, then say \emph{perek} again
\end{itemize}

\begin{tcolorbox}[breakable,colback=teal!4,colframe=teal!35!black,title={Before you move on}]
What does \bn{পেরেক} mean? (\textquotedblleft{}A nail\textquotedblright{}.) Say it once more.
\end{tcolorbox}
\section[\texorpdfstring{\bn{হাতুড়ি} (hātuṛi)}{হাতুড়ি (hātuṛi)}]{\bn{হাতুড়ি} (hātuṛi) — a hammer}

\label{lesson:BN-C111-haturi}



\begin{quote}
Before the new one: say the Bengali for a rope, then the Bengali for a nail.
\end{quote}

\subsection*{You'll want to know: \bn{হাতুড়ি}}
\textbf{\bn{হাতুড়ি}} — \emph{hātuṛi} — \textquotedblleft{}a hammer\textquotedblright{}.

\begin{cousinweb}
\textbf{\bn{হাতুড়ি}} (\emph{hātuṛi}) is a tadbhava, a word inherited and reshaped by sound change. It comes from \emph{*hastakūṭa}, \textquotedblleft{}hand hammer\textquotedblright{}, a form CDIAL reconstructs from its descendants (CDIAL 14028).
\end{cousinweb}

\begin{grammarlens}[title={What you've built}]
One more word for this chapter.
\end{grammarlens}

\subsection*{Your turn}
\begin{itemize}
\raggedright
  \item \emph{Say it:} \emph{hātuṛi}
  \item \emph{Say it:} \emph{hātuṛi}, once more
  \item \emph{Say it:} say \emph{hātuṛi}
  \item \emph{Recall:} say the Bengali for a rope, then the Bengali for a nail, then say \emph{hātuṛi} again
\end{itemize}

\begin{tcolorbox}[breakable,colback=teal!4,colframe=teal!35!black,title={Before you move on}]
What does \bn{হাতুড়ি} mean? (\textquotedblleft{}A hammer\textquotedblright{}.) Say it once more.
\end{tcolorbox}
\section[\texorpdfstring{\bn{কুড়ুল} (kuṛul)}{কুড়ুল (kuṛul)}]{\bn{কুড়ুল} (kuṛul) — an axe}

\label{lesson:BN-C111-kurul}



\begin{quote}
Before the new one: say the Bengali for a nail, then the Bengali for a hammer.
\end{quote}

\subsection*{You'll want to know: \bn{কুড়ুল}}
\textbf{\bn{কুড়ুল}} — \emph{kuṛul} — \textquotedblleft{}an axe\textquotedblright{}.

\begin{cousinweb}
\textbf{\bn{কুড়ুল}} (\emph{kuṛul}) is a tadbhava, a word inherited from Sanskrit and reshaped by sound change. It comes from \emph{kuṭhāra}, \textquotedblleft{}axe\textquotedblright{} (CDIAL 3244). CDIAL thinks the word probably Dravidian.
\end{cousinweb}

\begin{grammarlens}[title={What you've built}]
One more word for this chapter.
\end{grammarlens}

\subsection*{Your turn}
\begin{itemize}
\raggedright
  \item \emph{Say it:} \emph{kuṛul}
  \item \emph{Say it:} \emph{kuṛul}, once more
  \item \emph{Say it:} say \emph{kuṛul}
  \item \emph{Recall:} say the Bengali for a nail, then the Bengali for a hammer, then say \emph{kuṛul} again
\end{itemize}

\begin{tcolorbox}[breakable,colback=teal!4,colframe=teal!35!black,title={Before you move on}]
What does \bn{কুড়ুল} mean? (\textquotedblleft{}An axe\textquotedblright{}.) Say it once more.
\end{tcolorbox}
\section[\texorpdfstring{\bn{ছাতা} (chhātā)}{ছাতা (chhātā)}]{\bn{ছাতা} (chhātā) — an umbrella}

\label{lesson:BN-C111-chhata}



\begin{quote}
Before the new one: say the Bengali for a hammer, then the Bengali for an axe.
\end{quote}

\subsection*{You'll want to know: \bn{ছাতা}}
\textbf{\bn{ছাতা}} — \emph{chhātā} — \textquotedblleft{}an umbrella\textquotedblright{}.

\begin{cousinweb}
\textbf{\bn{ছাতা}} (\emph{chhātā}) is a tadbhava, a word inherited from Sanskrit and reshaped by sound change. It comes from \emph{chattra}, \textquotedblleft{}parasol\textquotedblright{} (CDIAL 4972).
\end{cousinweb}

\begin{grammarlens}[title={What you've built}]
One more word for this chapter.
\end{grammarlens}

\subsection*{Your turn}
\begin{itemize}
\raggedright
  \item \emph{Say it:} \emph{chhātā}
  \item \emph{Say it:} \emph{chhātā}, once more
  \item \emph{Say it:} say \emph{chhātā}
  \item \emph{Recall:} say the Bengali for a hammer, then the Bengali for an axe, then say \emph{chhātā} again
\end{itemize}

\begin{tcolorbox}[breakable,colback=teal!4,colframe=teal!35!black,title={Before you move on}]
What does \bn{ছাতা} mean? (\textquotedblleft{}An umbrella\textquotedblright{}.) Say it once more.
\end{tcolorbox}
